% covers_5languages.tex -- One-page visual demonstration of the CoppeTeX 4.0
% multilingual model: the five built-in main languages side by side, all
% sharing the same Portuguese institutional template, differing only in
% the title (and the body language, not visible on the cover).
%
% This file uses the article class deliberately -- it is a meta-document
% about the coppe class, not a thesis. It expects the five cover PNGs
% cover_pt.png, cover_en.png, cover_es.png, cover_fr.png,
% cover_it.png to live in the same directory (generated by pdftoppm
% from the example_<lang>.pdf demos).
\documentclass[a4paper,landscape,10pt]{article}
\usepackage[margin=1cm]{geometry}
\usepackage[T1]{fontenc}
\usepackage[utf8]{inputenc}
\usepackage{graphicx}
\usepackage{caption}
\captionsetup{font=small,labelfont=bf}
\pagestyle{empty}
\begin{document}
\begin{center}
  {\Large\bfseries CoppeTeX 4.0 -- modelo multilíngue\par}
  \vspace{1mm}
  {\small As cinco línguas principais embutidas, com a mesma identidade institucional UFRJ/COPPE.\par}
\end{center}
\vspace{2mm}
\noindent
\begin{minipage}[t]{0.195\linewidth}\centering
  \includegraphics[width=\linewidth]{cover_pt.png}\\[1mm]
  \textbf{brazilian}\\[-0.2em]
  {\scriptsize \texttt{\textbackslash documentclass[dsc]\{coppe\}}}
\end{minipage}\hfill
\begin{minipage}[t]{0.195\linewidth}\centering
  \includegraphics[width=\linewidth]{cover_en.png}\\[1mm]
  \textbf{english}\\[-0.2em]
  {\scriptsize \texttt{[english,dsc]}}
\end{minipage}\hfill
\begin{minipage}[t]{0.195\linewidth}\centering
  \includegraphics[width=\linewidth]{cover_es.png}\\[1mm]
  \textbf{spanish}\\[-0.2em]
  {\scriptsize \texttt{[spanish,dsc]}}
\end{minipage}\hfill
\begin{minipage}[t]{0.195\linewidth}\centering
  \includegraphics[width=\linewidth]{cover_fr.png}\\[1mm]
  \textbf{french}\\[-0.2em]
  {\scriptsize \texttt{[french,dsc]}}
\end{minipage}\hfill
\begin{minipage}[t]{0.195\linewidth}\centering
  \includegraphics[width=\linewidth]{cover_it.png}\\[1mm]
  \textbf{italian}\\[-0.2em]
  {\scriptsize \texttt{[italian,dsc]}}
\end{minipage}

\vfill
\noindent\footnotesize
\textbf{Observe que:} (1) o bloco institucional no topo
(\emph{Universidade Federal do Rio de Janeiro}, \emph{Instituto Alberto
Luiz Coimbra de Pós-Graduação e Pesquisa de Engenharia (COPPE)},
\emph{Programa de Engenharia de Sistemas e Computação}) permanece em
português em todas as capas, por se tratar de identidade institucional
da UFRJ;
(2) o título do trabalho acompanha o idioma principal escolhido na opção
da classe;
(3) a cidade e o ano no rodapé também ficam em português.

\hfill\textit{Imagens geradas a partir dos arquivos \texttt{example\_\{pt,en,es,fr,it\}.tex} em \texttt{dist/}.}
\end{document}
